\documentclass[../../../main.tex]{subfiles}

\begin{document}

\chapter{[N] Miscellaneous Goodies}

This note discusses some of the miscellaneous Number Theory topics
that are not necessarily the major part in olympiad NT, but would be
very useful when solving problems. At least for KMO, I am aware that
it is generally required to prove many of the contents here during
the write-ups, but some problems are virtually impossible to solve
without previous knowledge on these topics!

\section{\texorpdfstring{$p$}{p}-adic Valuation \& LTE Lemma}

LTE Lemma, which stands for Lifting the Exponent Lemma, concerns about
$p$-adic valuation of two integers. Before we get into the lemma,
let's discuss about $p$-adic valuation.

\begin{definition}{}{}
The $\mathbf{p}$\textbf{-adic valuation} of an integer $n$ and prime
$p$, denoted as $\nu_p (n)$, is the largest possible integer $a$ such
that $p^a \mid n$.
\end{definition}

Some examples could be $\nu_2 (48) = 3$ and $\nu_3 (54) = 3$. It is
worth noting that some sources write this as $V_p (n)$, so they most
likely refer to the same thing.

\begin{important}
A very famous example would be Legendre's formula which states the
following equation for positive $n$ and prime $p$.
\vspace{-0.8\baselineskip}
\[
\nu_p (n!)
= \sum_{k=1}^\infty \left\lfloor \frac{n}{p^k} \right\rfloor
\]
\end{important}

This is actually the same as the classic competition style problems
that asks you the number of $p$ in $n!$. I will skip the proof because
I think it is pretty self-evident.

Another very famous application is understanding why $\binom{n}{r}$ is
an integer. Obviously combinatoric-wise, its definition makes it an
integer, but algebraic-wise, it is pretty interesting considering the
fact that $\binom{n}{r} = \frac{n!}{r! (n-r)!}$. To show this,
it suffices to prove the following inequality for any prime $p$.
\[
\nu_p (n!) \geq \nu_p (r!) + \nu_p ((n-r)!)
\]
By Legendre's formula, the inequality converts as
\[
\sum_{k=1}^\infty \left\lfloor \frac{n}{p^k} \right\rfloor
\geq \sum_{k=1}^\infty \left(
    \left\lfloor \frac{r}{p^k} \right\rfloor +
    \left\lfloor \frac{n-r}{p^k} \right\rfloor
\right)
\]
and this is true because $\lfloor a + b \rfloor \geq \lfloor a \rfloor
+ \lfloor b \rfloor$. With this notation and concept in mind,
lets get into LTE Lemma!



\end{document}


